\section{Supplementary Material}
\label{sec:supp}

\subsection{A. Continuous Integration Matrix}
Table~\ref{tab:ci} summarizes the CI configuration that validates reproducibility
across Python versions. Each cell denotes a workflow executed and asserted against
the output contract.

\begin{table}[h]
\centering
\caption{Continuous-integration matrix. All three workflows are executed and their
output contract (success, report\_path, tables, figures) is asserted on every
Python version under \texttt{ubuntu-latest}.}
\label{tab:ci}
\begin{tabular}{lccc}
\toprule
Python & Literature & Biomarker & Causal \\
\midrule
3.9  & yes & yes & yes \\
3.10 & yes & yes & yes \\
3.11 & yes & yes & yes \\
3.12 & yes & yes & yes \\
\bottomrule
\end{tabular}
\end{table}

\subsection{B. Reproducibility Check}
The causal workflow accepts \texttt{--seed} (default 42). The \texttt{determinism}
CI job runs the causal workflow twice with \texttt{--seed 42} and asserts that the
two runs produce an identical output-file structure (same file paths and table
shapes). The literature workflow (\texttt{--use-mock}) and biomarker workflow
(\texttt{--use-synthetic}) run on fixed local fixtures and require no
network access.

\subsection{C. Output Categorization Schema}
Every extracted claim is assigned a five-level evidence grade using a rule-based scheme:

\begin{table}[h]
\centering
\caption{Output-category definitions used by the categorization scheme.}
\label{tab:evidence}
\begin{tabular}{lp{0.68\linewidth}}
\toprule
Grade & Definition \\
\midrule
A & Multiple independent high-quality studies, or direct experimental validation. \\
B & Single high-quality study, or strong consensus across sources. \\
C & Moderate evidence; plausible but with limited replication. \\
D & Weak or indirect evidence; hypothesis-level only. \\
E & Assertion unsupported by retrieved evidence; excluded from conclusions. \\
\bottomrule
\end{tabular}
\end{table}

\subsection{D. Twelve-Stage Execution Diagram}
The workflow runner executes twelve discrete stages: (1)~question formulation,
(2)~hypothesis generation, (3)~literature retrieval, (4)~study design,
(5)~data ingestion, (6)~statistical analysis, (7)~bioinformatics,
(8)~output aggregation, (9)~writing, (10)~internal review, (11)~critique,
(12)~formatted output. Each stage emits an execution output and passes through a validation gate.

\subsection{E. Standardized Interface Schema}
The bundled skill package exposes a single entry point (\texttt{run}) with a
model-agnostic tool schema. It returns the same output contract as the CLI and SDK
(Section~\ref{sec:repro}) and introduces no new computational capability. The
manifest pins \texttt{release\_tag = v1.0.0-arxiv-public} and
\texttt{frozen\_commit = cc5911e}.

\subsection{F. Provenance and Commit Binding}
This submission is a frozen snapshot:
\begin{itemize}
\item Core framework: commit \texttt{\corecommit} (branch \texttt{main}).
\item Standardized interface: commit \texttt{\skillcommit} (branch \texttt{feat/skills-layer}).
\item Paper source and bundled artifacts: release tag \texttt{\releasetag}
      (branch \texttt{\papersub}).
\end{itemize}
All artifacts are archived at the cited commits and are verifiable by reviewers.
